\chapter{Introduction}
\label{sec:introduction}

Large language models (LLMs) have become capable agents that can plan, reason, and call external tools to act in the world~\citep{brown2020gpt3,codeact2024}. The standard framing of this capability is \emph{tool use}: given a fixed set of pre-built functions, an agent selects the right one and passes the appropriate arguments~\citep{schick2023toolformer}. Benchmark suites such as MCP-Universe~\citep{mcpuniverse2025}, MCPAgentBench~\citep{mcpagentbench2024}, and WAPIIBench~\citep{wapiibench2025} measure this ability extensively, and the field has made substantial progress on it.

What these benchmarks leave unexamined is the step that must come before tool use in any realistic deployment: someone has to create the tools in the first place. In practice this is done by human engineers who read API documentation, write wrapper code, and expose the result through a protocol the agent can call. This process does not scale. A modern enterprise ecosystem spans hundreds of Representational State Transfer (REST) APIs; keeping tool wrappers current as those APIs evolve is a continuous maintenance burden. The gap between the APIs that exist and the APIs that agents can actually reach is, at present, almost entirely a human-labour bottleneck.

An appealing alternative is \emph{tool synthesis}: letting the model itself read documentation and produce the tool implementation. Contemporary models can write code~\citep{toolfactory2025}, operate as autonomous agents with file-system access~\citep{codeact2024}, and generate callable tool implementations from unstructured documentation~\citep{doc2agent2025}. Critically, these works generate individual tool stubs or callable functions; none produce a complete, deployable server exposing a typed protocol that an agent can discover and invoke at runtime. The combination suggests that producing a functional, complete MCP server from raw API documentation may already be within reach, yet no prior work has systematically evaluated this capability across diverse APIs under controlled conditions.

This shift from tool use to tool synthesis is conceptually framed by the \emph{LLM as Tool Maker} (LATM) paradigm~\citep{cai2023latm}, in which a capable model acts as an architect that proposes, verifies, and packages software tools for downstream use. The adoption of the Model Context Protocol (MCP;~\citeauthor{anthropic2024mcp},~\citeyear{anthropic2024mcp}) makes this practically tractable. Unlike traditional REST integrations, which require bespoke human-engineered connectors for each client--API pair, MCP wraps heterogeneous backends in a uniform JSON-RPC interface that agents can discover and call at runtime. This collapses integration complexity from $N \times M$ (every model client wired to every API) to $N + M$ (each side adapts once to the shared protocol), making automated synthesis the natural path to scaling agentic tool coverage.

A fundamental tension arises, however, at the boundary of neural generation and software execution. API interfaces are rigid: parameter names must match exactly, types must be correct, and HTTP payloads must conform to schema. LLMs generate code and serialise arguments through probabilistic token selection, which does not inherently respect these hard constraints. More importantly, the large presence of standardised API patterns and open-source library code in pre-training corpora creates strong, low-entropy priors over common software utilities. For well-known APIs such as GitHub or Stripe, a model may produce a functionally correct wrapper entirely from memory, ignoring the documentation provided at inference time. This phenomenon --- sometimes referred to as \emph{parametric gravity}~\citep{mallen2023knowns,juice2025} --- means that in-context documentation does not automatically override training-time knowledge, particularly for APIs that are heavily represented in pre-training data.

These observations motivate a systematic study. This thesis asks: \emph{can language models reliably synthesize functional API tool servers from documentation, and do those servers enable effective agent task completion?} To answer this, we introduce an end-to-end benchmark covering 14 REST APIs drawn from productivity tools, e-commerce platforms, communication systems, financial services, and developer tooling. The APIs are chosen to vary along two axes: \emph{familiarity} (the degree to which the API is likely represented in pre-training data) and \emph{complexity} (number of endpoints, parameter density, and structural depth). For each API, a synthesis agent is given raw documentation and asked to produce a complete, runnable MCP server. The resulting server is then
evaluated by a separate agent that attempts 13--25 natural-language tasks against it, with outcomes systematically judged by a third language model.

Formally, tool synthesis is an autoregressive generation process conditioned on a task specification $T$, the model's static weights $\theta$, and an external documentation context $C$:
%
\begin{equation}
  P(X \mid T, C;\, \theta)
  = \prod_{i=1}^{n} P(x_i \mid x_{<i},\, T,\, C;\, \theta),
  \label{eq:synthesis}
\end{equation}
%
where $X = (x_1, \ldots, x_n)$ is the generated server implementation. We study three instantiations of the context $C$: the full documentation $C_{\text{docs}}$, the empty context $C_{\emptyset}$ (no documentation), and a counter-factually mutated context $C_{\text{mut}}$ in which core parameter
names are renamed to conflict with training-time knowledge. Holding $T$ and $\theta$ fixed while varying $C$ isolates how much the output distribution depends on in-context evidence versus the parametric knowledge encoded in $\theta$.

Existing software engineering benchmarks such as SWE-bench~\citep{jimenez2023swebench} evaluate agents as black boxes, measuring macro pass/fail rates without explaining where and why failures occur. Conversely, mechanistic work on parametric memory and context--memory conflict~\citep{mallen2023knowns,longpre2021entity,juice2025} is largely restricted to open-domain question-answering datasets, which do not capture the stateful, protocol-driven constraints of live API execution. This thesis bridges that gap empirically: we measure parametric influence through controlled documentation ablations and counterfactual mutations, using a live execution environment as the ground truth signal. A preliminary logit-lens probing study on a smaller open-weight model localises parametric gravity to early-to-mid transformer layers, providing a mechanistic grounding for these behavioural patterns.

Based on these observations, the present thesis is essentially structured around five research questions:

\begin{itemize}
  \item \textbf{RQ1:} Can LLMs synthesize functional API tool servers from documentation alone?
  \item \textbf{RQ2:} Do synthesized servers enable effective agent task completion, compared to human-authored references?
  \item \textbf{RQ3:} To what extent do synthesizers rely on parametric knowledge rather than documentation, and does this affect performance?
  \item \textbf{RQ4:} How does documentation availability affect synthesis and agent performance?
  \item \textbf{RQ5:} Does synthesis quality predict agent task performance?
\end{itemize}

We also distil the findings into practical design guidelines for full-server MCP synthesis pipelines, covering documentation pre-processing, context inclusion decisions, and synthesis quality as a cheap pre-flight quality signal.

We report two primary metrics. \textbf{PASS} (agent pass rate) is the fraction of benchmark task attempts that the evaluation agent completes successfully, as judged by a separate LLM. \textbf{SYNTH} (synthesis sufficiency) is the fraction of tasks for which an LLM judge finds the synthesized server's tool list sufficient to complete the task, assessed without running the evaluation agent; it measures synthesis quality independently of agent reasoning ability.

Our main findings are: (i)~synthesis is feasible for most APIs under the best conditions (frontier-tier median PASS~83\%), and synthesized servers match or exceed human-authored reference implementations on GitHub and Zulip; (ii)~removing documentation does not degrade---and sometimes improves---PASS on high-frequency APIs, because model weights already encode correct implementations and documentation adds context noise; (iii)~mutation adherence is near-zero for well-known public APIs, confirming that parametric knowledge dominates documentation for those APIs; (iv)~SYNTH and PASS are strongly correlated ($r = 0.937$, $n = 209$ across all models), establishing the LLM verdict on synthesis quality as a reliable and cheap proxy for downstream agent performance. We additionally present a preliminary probing analysis connecting parametric dominance to early-to-mid transformer layers, and derive practical guidelines for robust synthesis system design.

The remainder of this thesis is organised as follows.
Chapter~\ref{sec:related_work} surveys related work on tool-use benchmarks, code generation, parametric knowledge, and the MCP protocol.
Chapter~\ref{sec:methodology} describes the benchmark construction, synthesis pipeline, evaluation harness, and experimental conditions.
Chapter~\ref{sec:results} presents results structured around the five research questions.
Chapter~\ref{sec:mechanistic} provides a logit-lens probing study localising parametric gravity to specific transformer layers.
Chapter~\ref{sec:discussion} interprets the findings in light of prior work, discusses practical consequences, and reflects on limitations.
Chapter~\ref{sec:conclusion} answers each research question, states the contribution, and identifies directions for future work.

